\section{Setting and Notation}
\label{sec:setting}

\paragraph{Field, cell ring, and trace.} Fix a prime $q$ with $q \gg 2^{32}$
and write $\F_q$ for the field of size $q$. Fix a row count
$n = 2^{\nu}$. The \emph{cell ring} is
\[
  \cellring \;:=\; \F_2[X]\big/(X^{32}),
\]
a free $\F_2$-module of rank $32$ and an $\F_2$-algebra. Its
$\F_q$-lift is
\[
  \liftring \;:=\; \cellring \otimes_{\F_2} \F_q \;\cong\;
  \F_q[X]\big/(X^{32}),
\]
again a free $\F_q$-module of rank $32$. A \emph{trace column} is an
element $v \in \cellring^n$, indexed by row $i \in \{0,1\}^{\nu}$.

\paragraph{Multilinear extension.} For $v \in \cellring^n$, the
\emph{multilinear extension} of $v$ is the unique polynomial
$\mle{v} \in \liftring[X_1,\ldots,X_{\nu}]$ that is multilinear in
$(X_1,\ldots,X_{\nu})$ and agrees with $v$ on $\{0,1\}^{\nu}$:
\[
  \mle{v}(r) \;=\; \sum_{i \in \{0,1\}^{\nu}} \eq_i(r) \cdot v_i,
  \qquad \eq_i(r) := \prod_{k=1}^{\nu} \bigl(i_k r_k + (1-i_k)(1-r_k)\bigr).
\]
The coefficients $\eq_i(r) \in \F_q$ depend only on $r$ and the row
index $i$; the cell value $v_i \in \cellring$ is $\F_2$-embedded in
$\liftring$. We use the same notation $\mle{v}(r)$ for the value at a
challenge $r \in \F_q^{\nu}$, which is an element of $\liftring$.

\paragraph{$\psia$-projection.} Fix a transcript-bound element
$\alpha \in \F_q$. Define the $\F_q$-linear projection
\[
  \psia : \liftring \;\longrightarrow\; \F_q,
  \qquad
  \psia\!\Big(\textstyle\sum_{j=0}^{31} c_j X^j\Big) \;=\; \sum_{j=0}^{31} c_j \alpha^j.
\]
Composed with $\F_2 \hookrightarrow \F_q$, it descends to a map
$\cellring \to \F_q$ which we still denote $\psia$. The PIOP commits
the multilinear extension of $\psia(v)$, not of $v$ itself.

\paragraph{Lifted evaluation.} For $v \in \cellring^n$ and
$r \in \F_q^{\nu}$, define the \emph{lifted evaluation}
\[
  \lmle{v}(r) \;:=\; \mle{v}(r) \;\in\; \liftring.
\]
At the end of the PIOP, the prover opens $\lmle{v}(r_0)$ for each base
column $v$ at a final point $r_0 \in \F_q^{\nu}$, providing all $32$
$\F_q$-coefficients of the lifted evaluation. The verifier checks that
$\psia(\lmle{v}(r_0))$ equals the value committed by the multi-point
evaluation argument.

\paragraph{Constraints and trace rows.} A UAIR specifies a constraint
polynomial $C(\mathrm{up}, \mathrm{down})$ where $\mathrm{up}$ and
$\mathrm{down}$ are abstract trace rows. The $\mathrm{up}$ row exposes
the cells of trace columns at row $i$; the $\mathrm{down}$ row exposes
\emph{virtual} cells, namely:
\begin{itemize}[leftmargin=2em]
  \item \emph{Row-shift virtual columns:} $\rowsh{\delta} v$ at row
    $i$ is $v_{i+\delta}$. These are already part of the protocol.
  \item \emph{Bit-op virtual columns:} the construction of this paper.
\end{itemize}

\paragraph{Index sets.} Let
$\Sup = \{1,\ldots,K\}$ be the set of base trace column indices.
Let $\Ssh = \{(j,\delta_\ell)\}$ be the set of row-shift specs. We
introduce a third index set
\[
  \Sbit \;=\; \{ \bitspec{j_\ell}{\bitop_\ell} \}_{\ell=1}^{m}
\]
of \emph{bit-op specs}, where $j_\ell \in \Sup$ is a base column index
and $\bitop_\ell$ is a bit-op (defined next).
